%=============================================================
%  General Conclusion
%=============================================================
\chapter*{General Conclusion}
\addcontentsline{toc}{chapter}{General Conclusion}

\section*{Summary of Work}

This project set out to design and implement \textbf{LUMINA}, a full-stack
intelligent e-commerce platform targeting the Tunisian cosmetics market ---
a sector where no locally-tailored, AI-augmented solution currently exists.

Over three Scrum iterations, the following was delivered:

\begin{itemize}
  \item \textbf{Iteration 1 --- E-Commerce Core}: a complete shopping experience
        including email OTP registration, TOTP two-factor authentication,
        a fully filtered product catalogue backed by PostgreSQL JSONB queries,
        a persistent cross-device cart, an atomic checkout transaction, and
        an admin product management interface.

  \item \textbf{Iteration 2 --- AI \& Personalisation}: a RAG chatbot grounded
        in the live product catalogue, protected by a three-layer prompt
        injection guard (length, regex, ML classifier), paired with a
        behavioral signal recommendation engine that personalizes the
        product feed without any external machine learning service.

  \item \textbf{Iteration 3 --- Administration \& Security}: a full admin
        dashboard for user and order management, JWT token revocation via
        Redis blacklisting with fail-closed behavior, role-based access
        control enforced at the controller layer, and Zod input validation
        across all API endpoints.
\end{itemize}

The system is built on a Clean Architecture foundation that enforces strict
separation between domain logic, application use cases, and infrastructure
concerns --- making it straightforward to swap out any technology layer
(ORM, LLM, payment provider) without touching business rules.

\section*{Technical Contributions}

Beyond the product itself, this project produced several reusable patterns:

\begin{itemize}
  \item An \textbf{abstract \texttt{ITransactionManager}} that allows use cases
        to run atomic database transactions without importing any ORM ---
        a clean-architecture-compliant solution to a common problem.
  \item A \textbf{three-layer prompt injection guard} applicable to any
        LLM-facing endpoint, combining lightweight and deterministic checks
        before invoking the costlier ML classifier.
  \item A \textbf{fail-closed JWT blacklist} pattern using Redis TTLs aligned
        to token expiration, ensuring zero windows of unauthorized access
        even in the event of infrastructure failure.
\end{itemize}

\section*{Limitations}

Honest assessment of the current state of the platform surfaces several
limitations that were consciously deferred:

\begin{itemize}
  \item \textbf{Card payments} --- Stripe integration is stubbed and explicitly
        rejects card payment requests. The infrastructure (PaymentIntent flow,
        webhook verification) is designed but not wired, pending PCI DSS review.
  \item \textbf{LLM latency} --- Local Ollama inference on CPU introduces
        3--8 second response times for the chatbot. A GPU-equipped server
        or a migration to a hosted API (OpenAI, Mistral) would resolve this
        in production.
  \item \textbf{Recommendation depth} --- The current engine uses a weighted
        signal sum. A collaborative filtering or matrix factorization model
        would improve cold-start performance for new users.
  \item \textbf{Deployment} --- The platform is fully containerized with Docker
        but has not yet been deployed to a public cloud environment.
\end{itemize}

\section*{Perspectives}

The following improvements are planned for future versions of LUMINA:

\begin{itemize}
  \item \textbf{Stripe payment integration} --- completing the PaymentIntents
        flow with webhook-based order confirmation and refund support.
  \item \textbf{Collaborative filtering} --- replacing the signal-weight model
        with a user-item matrix factorization approach (ALS or SVD)
        for better personalization at scale.
  \item \textbf{Real-time notifications} --- WebSocket or SSE-based order status
        updates and stock alerts for customers and admins.
  \item \textbf{Mobile application} --- a React Native client sharing the
        same API layer, extending LUMINA's reach to mobile-first users
        in a market where over 70\% of browsing happens on smartphones.
  \item \textbf{Analytics pipeline} --- an event streaming layer (Kafka or
        Redis Streams) feeding a business intelligence dashboard with
        cohort analysis, funnel metrics, and revenue forecasting.
\end{itemize}

\section*{Personal Reflection}

This project was as much an exercise in software architecture as it was in
product thinking. Designing a system that is simultaneously maintainable,
secure, and intelligent --- while staying within the constraints of local
infrastructure and academic timelines --- required constant prioritization
and trade-off analysis.
The Scrum methodology proved invaluable in this regard: short iterations
forced early delivery of working software and surfaced integration issues
before they became blockers.

\vspace{1em}
\noindent
\textbf{LUMINA} demonstrates that building a production-quality, AI-augmented
e-commerce platform is achievable without expensive cloud AI subscriptions,
without compromising user data privacy, and without sacrificing the
engineering discipline that makes software sustainable in the long run.
